\section{Institutional Background}
\label{sec:background}

\subsection{Tour Structure}

The Association of Tennis Professionals (ATP) organizes men's professional tennis and the Women's Tennis Association (WTA) organizes women's professional tennis, each into a hierarchical tournament structure. Four Grand Slam tournaments (Australian Open, Roland Garros, Wimbledon, US Open) sit at the top, followed by Masters 1000 events, and then lower-tier events (ATP 500/250; WTA 500/250). Each tournament awards ranking points based on the round reached, with Grand Slams offering the most (2,000 for the winner) and lower-tier events offering fewer.

The two tours differ in ways that matter for the analysis. The WTA calendar is denser, with more tournaments per month and shorter gaps between events. WTA qualifying draws are typically smaller---often 12 players compared to 16 on the ATP side---which means fewer final-round qualifying losers per event. I return to these differences when interpreting the heterogeneous treatment dynamics across tours in Section~\ref{sec:results_heterogeneity}.

\subsection{Ranking Systems and the Access Feedback Loop}
\label{sec:background_rankings}

Both tours use a rolling 52-week ranking system. A player's ranking equals the sum of points earned across their best results in the previous 52 weeks. Rankings are updated weekly and determine direct entry into tournament main draws: the top 104 players enter Grand Slam main draws directly, the top 40--56 enter Masters events, and similar thresholds apply at lower tiers.

This structure creates a feedback loop that is central to the analysis. Ranking points determine tournament access, and tournament access determines opportunities to earn ranking points. A player who earns unexpected points at one event may gain direct entry to subsequent events they would otherwise have needed to qualify for, generating further point-earning opportunities. This form of cumulative advantage \citep{Merton1968_matthew_effect} means that a single exogenous shock to ranking points can propagate through the system---at least temporarily.

The feedback loop also has an ``overseeding'' property. A player whose ranking exceeds their ability is placed into draws with opponents who are, on average, stronger than what the player's skill warrants. The tournament performance model in Section~\ref{sec:mechanisms_performance} tests and confirms this prediction.

\subsection{Qualifying Draws}

Players ranked below the direct-entry threshold can enter a tournament's qualifying draw, a mini-tournament held in the days before the main event. Qualifying draws are typically single-elimination brackets (32 or 16 players on the ATP tour, 12 or 16 on the WTA tour), with 4--8 qualifiers advancing to the main draw. To qualify, a player must win two or three consecutive matches.

\subsection{The Lucky Loser Mechanism}
\label{sec:background_ll}

When a player who has been accepted into a tournament main draw withdraws---due to injury, illness, or scheduling conflicts---the vacant slot must be filled. The replacement comes from the pool of ``lucky losers'' (LLs): players who lost in the final round of qualifying.

\paragraph{Grand Slam tournaments.} The Grand Slam rulebook specifies a two-track system based on withdrawal timing. When a main draw withdrawal occurs \textit{before} qualifying is complete, the top four ranked final-round qualifying losers enter a random draw (lottery) for the available slot. Each of the four has an equal probability of selection regardless of ranking differences within the group. When a withdrawal occurs \textit{after} qualifying ends, the highest-ranked final-round qualifying loser receives the slot without a lottery. The lottery pathway provides random assignment and is the basis for the primary identification strategy. The same mechanism operates at both ATP and WTA Grand Slams. Lottery-based assignment commenced after Wimbledon 2005; the Grand Slam lottery sample therefore covers 2006--2024.

\paragraph{Non-Grand-Slam events.} At Masters 1000, ATP 500, and ATP 250 tournaments, lucky loser selection follows a strict ranking rule: the highest-ranked final-round qualifying loser receives the first available slot, the second-highest receives the next, and so on. There is no lottery component.

\subsection{Elo Ratings}
\label{sec:background_elo}

I use Elo ratings as a measure of playing ability that is independent of the ranking point system. Originally developed for chess \citep{Klaassen2001_tennis_iid}, the Elo system updates player ratings after each match based on the match outcome and the expected outcome given both players' pre-match ratings. A player gains rating points for winning and loses them for losing, with the magnitude depending on the opponent's strength. I compute Elo ratings from the full match history using a $K$-factor of 32 and an initial rating of 1,500.

The distinction between ranking points and Elo ratings is central to the mechanism analysis. Ranking points are a direct function of tournament participation and rounds won: a player who enters a main draw and wins one match earns a fixed number of points regardless of the opponent's strength. Elo ratings adjust for opponent quality: beating a higher-rated player produces a larger Elo gain than beating a lower-rated one. If lucky loser entry improves rankings but not Elo ratings, the improvement reflects more participation rather than improved ability to win matches against quality-adjusted opponents.
